\documentclass[10pt,a4paper]{amsart}
\usepackage[margin=19mm]{geometry}
\usepackage{amsmath,amssymb,mathtools}
\usepackage[scr=rsfs,cal=euler]{mathalfa}
\usepackage{xfrac}
\usepackage[unicode,hidelinks,bookmarks=false]{hyperref}
\newcommand{\ie}{i.e.\ }
\newcommand{\cf}{cf.\ }
\newcommand{\resp}{respectively\ }
\newcommand{\Subsection}{\subsection}
\providecommand{\nobreakdash}{\nobreak\hbox{-}}
\setlength{\emergencystretch}{2em}
\multlinegap0pt
\usepackage{mathtools}
\usepackage{amssymb}
\usepackage{enumerate}
\newtheorem{lemma}{Lemma}
\newtheorem{theorem}{Theorem}
\newtheorem{proposition}{Proposition}
\newcommand{\Assoc}{\mathrm{Assoc}}
\newcommand{\Dbarn}{\overline{\mathrm{Diff}}_\partial(D^n)}
\newcommand{\Diff}{\mathrm{Diff}}
\newcommand{\Diffn}{\mathrm{Diff}_\partial(D^n)}
\newcommand{\Homeo}{\mathrm{Homeo}}
\newcommand{\OO}{\mathrm{O}}
\newcommand{\TOP}{\mathrm{TOP}}
\newcommand{\bbslash}{\backslash\!\!\backslash}
\newcommand{\hofiber}{\mathrm{hofiber}}
\newcommand{\mmod}{\,\mathrm{mod}\,}
\newcommand{\pdHomeo}{\mathrm{pdHomeo}}
\newcommand{\plHomeo}{\mathrm{plHomeo}}
\title{On the smoothing theory delooping of disc diffeomorphism and embedding spaces}
\author{Paolo Salvatore, Victor Turchin}
\date{Source: arXiv:2407.14699v4 (CC BY 4.0)}
\begin{document}
\maketitle

\noindent\emph{Source and licence.} On the smoothing theory delooping of disc diffeomorphism and embedding spaces. Version arXiv:2407.14699v4 (CC BY 4.0), distributed by arXiv under the Creative Commons Attribution 4.0 International licence, http://creativecommons.org/licenses/by/4.0/, as recorded in the arXiv licence metadata for this exact version and shown on the abstract page https://arxiv.org/abs/2407.14699v4. Full licence notice: \texttt{licenses/arxiv-2407.14699-CC-BY-4.0.txt}.\par\smallskip

\noindent\textit{Editorial scope.} Counted unit: the article's theorem th:a\_top (source0.tex lines 1939--1947). The article's complete original proof of it is delivered with it (source0.tex lines 1949--1974, one \texttt{proof} environment of 26 lines). No mathematics is written, repaired or re-derived: the statement and the proof are the article's own lines, in the article's own notation, and no retained line is substituted or re-flowed.  Retained uncounted as the source-local context the proof needs: lines 735--740; lines 1073--1073; lines 1189--1189; lines 1246--1246; lines 1298--1298; lines 1773--1775; lines 1790--1792; lines 1814--1816; lines 3072--3079.  Nothing was omitted from any retained span, so the package carries no omission record.  No retained line contains a citation, so the wrapper carries no bibliography and no bibliography entry.  No retained line contains a figure environment, an image-inclusion command or drawing code, so no image is delivered and the package declares no asset dependency.  Editorial formatting only, in addition: an AMS \texttt{amsart} preamble that restates the article's own declarations and macros used by the retained lines, namely the environments \texttt{lemma}, \texttt{theorem}, \texttt{proposition} on separate counters, together with the standard packages those lines need.  The retained environments print in this document as Lemma 1, Theorem 1, Proposition 1 in order of appearance, whereas the article prints the same items with the numbering of its own full document.  The complete native source of the article as distributed in the arXiv e-print for this exact version is bundled unchanged as \texttt{sources/162606/source0.tex}.
\begin{lemma}[Alexander trick]\label{l:alex}
For any $0\leq m\leq n$, the groups $\Homeo_\partial(D^n)$, $\plHomeo_\partial(D^n)$, %$\pdHomeo_\partial(D^n)$,
 $\Homeo_\partial(D^n\mmod D^m)$, 
$\plHomeo_\partial(D^n\mmod D^m)$ %, $\pdHomeo_\partial(D^n\mmod D^m)$
 are contractible.
\end{lemma} 

\subsection{Framed operads and their action}\label{ss:fr_operads}

\subsection{Homotopy fibers and pullbacks}\label{ss:hofib}

\subsubsection{Loop spaces of  homotopy fibers and pullbacks}\label{sss:loops}

\subsection{Embeddings modulo immersions spaces}\label{ss:modulo} 

\begin{equation}\label{eq:dif_fr_homeo}
\Diff_\partial(N)\to \Homeo_\partial^{\underline{O_n}}(N).
\end{equation}

\begin{lemma}\label{l:dif_fr_homeo}
For any smooth manifold $N$ (with Riemannian metric), the map \eqref{eq:dif_fr_homeo} is a homomorphism of simplicial groups.
\end{lemma}

\begin{theorem}\label{th:diff_compos}
For any smooth manifold $N$ of dimension $n\neq 4$, the map $\Diff_\partial(N)\to \Homeo_\partial^{\underline{O_n}}(N)$  is an equivalence of simplicial groups.
\end{theorem}

\begin{proposition}\label{l:E1equiv}
Let $H,K\subset G$ be subgroups of a topological group~$G$, and let $N\subset N(H)\cap N(K)$ be a subgroup of the intersection of their normalizers. We additionally assume that all these (sub)groups are cofibrant as spaces
and all the underlying group inclusions are cofibrations of spaces.\footnote{For example, $G$ is a Lie group, while $H,K,N$ are its closed subgroups. As another example relevant to us, $G,H,K,N$ can be realizations of simplicial (sub)groups.} One then has an equivalence
of $E_1$-algebras in $N$-spaces:
%\begin{equation}\label{eq:E1equiv}
$H\times_G^h K\simeq_{E_1^N}\Omega\left(H\bbslash G/K\right)$.
%\end{equation}
\end{proposition}

\begin{theorem}\label{th:a_top}
For $n\neq 4$, one has an equivalence of $E_{n+1}^{\OO_n}$-algebras
\begin{equation}\label{eq:a_top_1}
\Diffn\simeq_{E_{n+1}^{\OO_n}}\Omega^{n+1}\left(\TOP_n/\OO_n\right);
\end{equation}
\begin{equation}\label{eq:a_top_2}
\Dbarn\simeq_{E_{n+1}^{\OO_n}}\Omega^{n+1}\TOP_n.
\end{equation}
\end{theorem}

\begin{proof}[Proof of~\eqref{eq:a_top_1}]
Consider the following zigzag
\[
\Diffn\xrightarrow[I]{\,\simeq\,}\Homeo_\partial^{\underline{O_n}}(D^n)\xleftarrow[J]{\,\simeq\,}\Omega^n\hofiber(O_n\to\TOP_n).
\]
The first map is an equivalence by Theorem~\ref{th:diff_compos}. One can see that it respects the $E_n$-action. By Lemma~\ref{l:dif_fr_homeo} it  also respects the associative product, implying that $I$ is an equivalence of 
$R_{n+1}$-algebras. It is also not hard to see that $I$ respects the $\OO_{n}$-action, where the $\OO_{n}$-action on $\Diffn$ is described in Subsection~\ref{ss:fr_operads}, while the $\OO_n$-action on 
$\Homeo_\partial^{\underline{O_n}}(D^n)$ is defined in the same manner as for $\Dbarn$,
see Subsection~\ref{ss:modulo}. The map $J$ is the fiber inclusion  for the projection
$\Homeo_\partial^{\underline{O_n}}(D^n)\to \Homeo_\partial(D^n)$, the fiber 
%(see Subsections~\ref{ss:hofib} and~\ref{sss:loops}) 
 being taken over the
identity homeomorphism $id\in \Homeo_\partial(D^n)$. Since the base space $\Homeo_\partial(D^n)$ is contractible by Lemma~\ref{l:alex}, the fiber inclusion~$J$ is an equivalence. On the other hand, $\{id\}$
is an $R_{n+1}^{\OO_n}$-subalgebra of $\Homeo_\partial(D^n)$ and the projection is a morphism
of $R_{n+1}^{\OO_n}$-algebras. Thus, $J$ is also an inclusion of $R_{n+1}^{\OO_n}$-algebras.

To get the last step, i.e., equivalence to $\Omega^{n+1}\left(\TOP_n/ \OO_n\right)$, we pass from
simplicial sets to topological spaces. %changed
 We therefore first replace $\Omega^n\hofiber\left(O_n\to\TOP_n\right)$
by $\Omega^n\hofiber\left(|\OO_n|\to|\TOP_n|\right)$, see Subsection~\ref{sss:loops}. The new space
is an $R_{n+1}^{|\OO_n|}$-algebra ($|\OO_n|$ acts on $R_{n+1}$ by restriction along the group homomorphism $|\OO_n|\to \OO_n$). On the other hand, $\hofiber\left(|\OO_n|\to|\TOP_n|\right)$ is an $\Assoc^{|\OO_n|}$-algebra
$E_1^{|\OO_n|}$-equivalent to 
$\Omega\left(|\TOP_n|/|\OO_n|\right)=\Omega\left|\TOP_n/\OO_n\right|$ by Proposition~\ref{l:E1equiv} (where $|\OO_n|$ acts trivially on $\Assoc$ and $E_1$). We conclude that 
 $\Omega^n\hofiber\left(|\OO_n|\to|\TOP_n|\right)$ is $E_{n+1}^{|\OO_n|}$-equivalent to 
$\Omega^{n+1}\left|\TOP_n/\OO_n\right|$.
\end{proof}

\end{document}
